\section{GNSS・衛星測位}

\subsection{平成27年度 問33}
\label{q:h27-reg-applied-33}

\noindent\textbf{原典：} sources/H27/17\_applied\_science.pdf、PDF物理13頁。

\textbf{公式正答：} 3（single）。

\paragraph{問題の要旨} 測位衛星の時刻信号で位置を求める衛星測位システム全般の略称を選ぶ。

\paragraph{問われる核心} 衛星測位の総称と個別システム

\paragraph{必要知識} GNSSとGPSの関係、電波到達時間と擬似距離、受信機時計誤差。

\paragraph{考え方（査読待ち初稿）} 問われているのは測位を実現する国別の装置群ではなく、それらを含む名称である。米国GPS、ロシアGLONASS、欧州Galileo、日本QZSSを区別すると、総称はGNSSである。

\paragraph{正答を選ぶ理由} 問われているのは測位を実現する国別の装置群ではなく、それらを含む名称である。米国GPS、ロシアGLONASS、欧州Galileo、日本QZSSを区別すると、総称はGNSSである。

\paragraph{選択肢の検討（査読待ち初稿）}\mbox{}\par

\noindent\textbf{1：不適切} GALILEO。欧州の個別システム名であり、全般の総称ではない。\par

\noindent\textbf{2：不適切} GLONASS。ロシアの個別システム名である。\par

\noindent\textbf{3：適切} GNSS。Global Navigation Satellite Systemの略で、衛星測位システムの総称に対応する。\par

\noindent\textbf{4：不適切} GPS。米国の個別システム名である。\par

\noindent\textbf{5：不適切} QZSS。日本の準天頂衛星システムである。GNSSに関連する地域的システムを含めた総称と、その一システムを区別する。\par

\paragraph{再利用する考え方} 衛星測位の総称と個別システムを、測定量・推定対象・基準面・近似条件に分けて理解する。名称だけで探査法を対応させない。

\paragraph{共通理論} \hyperref[theory:gravity:gnss]{第\ref*{theory:gravity:gnss}節}

\par 参考：\cite{gsi_gnss}。

\par 難易度 3（校準初稿）、優先度未評価。
